% =====================================================================
%  GMN3 附录 F  Fock--Goncharov 坐标与 traffic rules
%        附录 G  $SU(2)$、$N_f=0$ 与切比雪夫（Tchebyshev）多项式
%  原文：Lamination-Paper.tex 6896--7129 行；PDF p.115--124（印刷 p.114--123）
%  编号公式 (F.1)--(F.10)、(G.1)--(G.8) 共 18 条，全部落在本块内，
%  故块内引用一律用 \eqref / \ref，label 名照抄原文；
%  跨块引用（式 (6.3)(6.14)(9.11)(10.24)、图 29、第 10.3 节）暂写原文编号，
%  行尾留 % TODO 合稿时统一改 \eqref / \ref。
%  附录字母由 \appendix 自动给出（\appendix 在 ch15），本块两节即 F、G。
% =====================================================================

\section{Fock--Goncharov 坐标与 traffic rules}

在式 (6.3)% TODO 合稿改 \eqref{eq:trh-large-r}
中，我们把某些线算子 vev $\inprod{L_\zeta}$ 等同于平坦联络绕 $C$ 上闭回路之和乐的迹。在式 (6.14)% TODO 合稿改 \eqref{eq:y1}
中，我们类似地把 vev $\inprod{L_\zeta}'$ 等同于扭平坦联络之和乐的迹。本附录说明如何用相应的 Fock--Goncharov 型坐标来展开这些和乐。这一计算的基本想法并不新鲜，在许多地方都能找到；唯一可能算新颖的一点是：在扭的情形下，即使和乐取 $SL(2,\IC)$ 的基础表示，我们也能够定义正的展开。

\subsection{坐标}\label{app:fgcoords}

我们先简要回顾 Fock--Goncharov 坐标的定义。这些坐标由文献 \cite{MR2233852} 引入；一个略作改写、便于我们使用的版本在文献 \cite{Gaiotto:2009hg} 中作了描述。

设给定 $C$ 的一个「装饰三角剖分」$T$。这指的是 $C$ 的一个三角剖分，其中每个顶点或是一个正则奇点，或是挖去一个非正则奇点的小圆盘上的一个标记点。「装饰」指的是联系于每个顶点的一种离散选择。对正则奇点，我们选取单值性的两个本征空间之一；对非正则奇点，我们选取标记点与自该奇点发出的斯托克斯线之间的一个对应。无论哪种情形，给定一个顶点与一个平坦联络 $\CA$，装饰的选取都给出一种办法，挑出一个平坦截面 $s$（即 $(d+\CA)s=0$ 的解），允许差一个标量倍数。对正则奇点，$s$ 取为单值性的一个本征向量；对非正则奇点，它是沿反斯托克斯（又名 WKB）射线指数衰减的截面。由于单值性，$s$ 未必能在 $C$ 上整体存在，但至少可以取它在任何单连通区域上存在。

给定 $C$ 的任一装饰三角剖分 $T$ 与任一边 $E\in\Edges(T)$，我们考虑以 $E$ 为对角线的四边形 $Q_E$。从 $E$ 的某一端出发、按逆时针方向把四个顶点编号为 $1$ 到 $4$。于是 $T$ 的装饰给出四个平坦截面 $s_i$。我们可以把它们平行移动到 $Q_E$ 中的任一公共点，再计算它们 $SL(2,\IC)$ 不变的交叉比：
\begin{equation} \label{eq:def-Y}
\CY_E = \frac{(s_1 \wedge s_2) (s_3 \wedge s_4)}{(s_2 \wedge s_3)(s_4 \wedge s_1)}.
\end{equation}

对扭局部系统的情形，我们也想要上述东西的一个版本。这时平坦截面 $\ts_i$ 定义在 $C$ 的穿刺切丛上——更确切地说，定义在其某些单连通子集上。因此，对如何搬运平坦截面 $\ts_i$，我们必须更加小心。我们取从各顶点到 $Q_E$ 内部一个公共点 $\CP_*$ 的路径，要求它们不自交、彼此也不相交。这之后，我们考虑 $\CP_*$ 上方的切向量所成的圆：在这个圆上我们有一个秩 $2$ 的平坦联络，其和乐为 $-1$，另有四个向量 $\ts_i$，位于圆上循环排序的点 $t_i$ 处。我们要定义这四个向量的「扭交叉比」。我们采用如下约定：楔积 $\ts_i \wedge \ts_{i+1}$ 的定义方式是，沿 $t_i$ 与 $t_{i+1}$ 之间不含另外两个 $t_j$ 的那段弧（即绕圆走「近路」）把两个向量平行移动到公共点，取它们的楔积，再除以 $SL(2,\IC)$ 不变的体积形式。按此约定可写
\begin{equation} \label{eq:def-tY}
\tCY_E = \frac{(\ts_1 \wedge \ts_2) (\ts_3 \wedge \ts_4)}{(\ts_2 \wedge \ts_3)(\ts_4 \wedge \ts_1)}.
\end{equation}


\subsection{traffic rules}\label{app:traffic}

在文献 \cite{Gaiotto:2009hg} 的附录 A 中，我们回顾了一个众所周知的算法，它用联系于某个三角剖分的 Fock--Goncharov 坐标来计算和乐及其迹。我们称这个算法为「traffic rule 算法」。这里我们描述为覆盖扭局部系统与叶层这两种情形所需的 traffic rules 的推广。

固定一个三角剖分 $T$ 与一个扭局部系统，即 $\tCM$ 中一点。再固定一条边 $E$ 以及 $E$ 的一个余定向 $v$。对 $(E,v)$ 决定了单位切丛中一个单连通区域 $D(E,v)$，它丛化于 $E$ 之上（由指向 $E$ 某一侧的全部切向量组成；为简便起见，我们也用 $v$ 记其中任一切向量）。设 $l$ 与 $r$ 是 $E$ 的两个顶点，其取法使沿 $E$ 从 $l$ 走向 $r$ 时 $v$ 在左侧。我们关心扭局部系统在 $D(E,v)$ 上的平坦截面。它们构成二维空间 $S(E,v)$，其中有两个由顶点处装饰确定的区分向量 $\ts_l$ 与 $\ts_r$（相差整体标量，按惯例任意固定）。

于是可以定义量 $N(E) = \frac{\ts_l \wedge \ts_r}{\vol}$，其中 $\vol$ 是 $SL(2,\IC)$ 不变的体积形式。关键在于，如此定义的 $N(E)$ 不依赖于余定向 $v$ 的选取。确实，反转 $v$ 会交换 $\ts_l$ 与 $\ts_r$，从而带来一个负号；但反转 $v$ 也改变了我们计算楔积的区域，而如图 \ref{fig:twisted-sign} 所示，这一改变会带来第二个负号。（正是在这里，使用\emph{扭}局部系统帮了我们。）

\insfig{twisted-sign}{截面 $\ts_l$ 与 $\ts_r$ 最初定义在两个顶点附近区域上单位切丛的可缩区域中。要计算楔积 $\ts_l \wedge \ts_r$，必须把它们搬运到 $C$ 中的同一公共点 $x$ 以及 $x$ 处的同一切向量。图中给出两种做法，对应于边之余定向的两种选取：或者把两个截面都向北搬，或者都向南搬，沿图中标出的路径。灰色圆表示单位切丛的纤维，路径被提升到这些纤维上；灰色圆上的小点表示提升后的切向量。这两种定义 $\ts_l \wedge \ts_r$ 的方式相差一个因子 $-1$：要看出这一点，注意当我们绕四条路径所成的回路走一圈时，切向量绕纤维转了一圈，而我们局部系统沿纤维的和乐为 $-1 \in SL(2,\IC)$。}

定义
\begin{equation}
\tau(E) = \sqrt{N(E)}
\end{equation}
其中我们一次性选定两个可能平方根之一；对每个 $E$ 这都是一个独立的选择。作了这些选择之后，我们也得到了 $\sqrt{\tCY_E}$ 的一个定义：沿用上一节的记号，我们固定
\begin{equation}
\sqrt{\tCY_E} = \frac{\tau(E_{12}) \tau(E_{34})}{\tau(E_{23}) \tau(E_{41})}.
\end{equation}

现在考虑位于某个面 $F$ 上的一条边 $E$。$F$ 决定了 $E$ 的两个余定向，即指向 $F$ 内部的 $v^{\mathrm{inn}}$ 与指向 $F$ 外部的 $v^{\mathrm{out}}$。设 $l$ 与 $r$ 是 $E$ 的两个顶点，取法使沿 $E$ 从 $l$ 走向 $r$ 时 $v^{\mathrm{inn}}$ 在左侧；设 $E_l$ 与 $E_r$ 是 $F$ 的另外两条边，分别含 $l$ 与 $r$。现在定义平坦截面空间的两组基：
\begin{align}
 B^{\mathrm{inn}}(E,F) &= \left( \ts_l \frac{\tau(E_r)}{\tau(E_l) \tau(E)}, \ts_r \frac{\tau(E_l)}{\tau(E_r) \tau(E)} \right), \\
 B^{\mathrm{out}}(E,F) &= \left( \ts_r \frac{\tau(E_l)}{\tau(E_r) \tau(E)}, \ts_l \frac{\tau(E_r)}{\tau(E_l) \tau(E)} \right),
\end{align}
它们分别定义在 $D(E, v^{\mathrm{inn}})$ 与 $D(E, v^{\mathrm{out}})$ 上。归一化因子选得使这两组基的行列式都等于 $1$。

我们要研究沿 $C$ 上路径 $p$ 的平行移动。把 $p$ 切成小段，就可以把该平行移动具体地写成矩阵的乘积，这些矩阵相对于 $p$ 所遇到的各种 $E$、$F$ 的基 $B^{\mathrm{inn},\mathrm{out}}(E,F)$ 而言。

在面 $F$ 上从边 $E$ 走到边 $E'$ 的一段路径，提升到单位切丛后给出从 $D(E, v^{\mathrm{inn}})$ 到 $D(E', v'^{\mathrm{out}})$ 的一段路径。相应的平行移动（相对于基 $B^{\mathrm{inn}}(E,F)$ 与 $B^{\mathrm{out}}(E',F)$）取决于路径是右转还是左转：用 Pl\"ucker 关系作一个简短的直接计算即得
\begin{equation}
 L = \begin{pmatrix} 1 & 1 \\ 0 & 1 \end{pmatrix}, \qquad R = \begin{pmatrix} 1 & 0 \\ 1 & 1 \end{pmatrix}.
\end{equation}
类似地，沿一段穿过边 $E$、从面 $F$ 进入面 $F'$ 的路径所作的平行移动，由矩阵
\begin{equation}
 M_E = \begin{pmatrix} \sqrt{\tCY_E} & 0 \\ 0 & 1 / \sqrt{\tCY_E} \end{pmatrix}
\end{equation}
给出。也就是说，$B^{\mathrm{inn}}(E,F') = B^{\mathrm{out}}(E,F) M_E$。见图 \ref{fig:turns}。

\insfig{turns}{联系于各段路径的平行移动矩阵。}

定义 $\tCY_p = \prod_E \tCY_E$，其中乘积遍历 $p$ 所穿过的全部边。由 $\tCY_E$ 的定义可见，$\tCY_p$ 是若干 $N(E)$ 的乘积，而每个 $N(E)$ 都取\emph{偶}次幂。这意味着 $\sqrt{\tCY_p}$ 可以典范地定义，不依赖任何选择。扭局部系统沿路径 $p$ 的和乐，是 $p$ 所穿过的全部边上的矩阵 $R$、$L$ 与 $M_E$ 的乘积。因此该和乐等于 $1/\sqrt{\tCY_p}$ 乘上一个矩阵，其元素是 $\tCY_E$ 的多项式，系数为\emph{正}整数。特别地，这意味着和乐的\emph{迹}等于 $1/\sqrt{\tCY_p}$ 乘上 $\tCY_E$ 的一个多项式，系数为正整数。

我们也能计算对应于叶层的线算子的期望值。回忆一下，叶层除了闭曲线之外，还可以含有终止于边界弧上的开曲线，并带有整数权重。这要求我们把 traffic rules 稍作扩充，以规定在开端点处该怎么做。事实上，对单独一条开曲线我们给不出简单的规则：我们必须考虑曲线的组合，使得终止于每条边界弧上的总权重为零。不失一般性，可以只限于这种组合中最简单的一个，见图 \ref{fig:lamination-basic-unit}。

\insfig{lamination-basic-unit}{开曲线的基本组合：在每条边界弧上的净权重为零。}

通过同痕，我们总可以安排成图中出现的四个端点都落在同一条边界边上。这样做之后，我们直接把带权重 $-1$ 的那段曲线丢掉（我们将选取归一化的截面，使这一段贡献 $1$）。至于另外两段曲线，我们先在每段上取一个定向，再按图 \ref{fig:lamination-ends} 的规则与下式给它们配上行向量或列向量（行向量表示 $B(E, v^{\mathrm{inn}})$ 中的基元素，列向量表示 $B(E, v^{\mathrm{out}})^*$ 中的基元素）：
\begin{gather}
 B^L = \begin{pmatrix} 0 & 1 \end{pmatrix}, \quad B^R = \begin{pmatrix} 1 & 0 \end{pmatrix}, \\
 E^L = \begin{pmatrix} 1 \\ 0 \end{pmatrix}, \quad E^R = \begin{pmatrix} 0 \\ 1 \end{pmatrix}.
\end{gather}

\insfig{lamination-ends}{一条开边可以开始或结束的四种方式，以及平坦截面空间中相应的基向量。}

这些公式直接来自叶层 vev 的定义（回忆一下，该定义说的是：把路径两端顶点所联系的平坦截面平行移动到同一公共点，再取它们的楔积）：这里我们把归一化的截面从路径的终点搬运到起点。


\section{$SU(2)$、$N_f = 0$ 与切比雪夫（Tchebyshev）多项式} \label{app:chebyshev}

本附录证明公式 (10.24)% TODO 合稿改 \eqref{eq:chebyshev-form}
，它给出 $SU(2)$、$N_f = 0$ 理论中 $\inprod{L_{2n}}'$ 的展开。我们随意使用第 10.3 节% TODO 合稿改 \ref{subsec:SU2-Lams}
的记号。

正如正文中已经提到的，为省事起见，我们避开「扭局部系统」这个词，改在 $\CM$ 而不是 $\tCM$ 上工作。代价是某些符号与平方根的选取会变得难以确定；我们干脆在最后以正性为要求把它们定下来，因为我们已经证明：在 $\tCM$ 上所需的展开确实具有正的系数。

把图 29% TODO 合稿改 \ref{fig:nf0-bare}
映射到一个环带上，于是有内圆与外圆。引入小的平坦截面 $s_1$（附着于外圆上的 WKB 射线）与 $s_2$（附着于内圆上的 WKB 射线）。以 $M$ 记顺时针的单值算子。于是把各因子 (9.11)% TODO 合稿改 \eqref{eq:open-vev}
组合起来得到
\begin{equation}
\langle L_{2n} \rangle = - \frac{ (s_2 \wedge M^n s_1)^2 }{(s_1 \wedge M s_1 )(s_2 \wedge M s_2) }.
\end{equation}
特别地，$\langle L_0 \rangle = 1/Y$，$\langle L_2 \rangle = X$。由此推出 $\inprod{L_1} = \sqrt{\frac{X}{Y}}$。

再引入 $M$ 的本征向量 $s_\pm$，使 $M s_\pm = \lambda^{\pm 1} s_{\pm}$，其中 $\lambda \in \IC^*$。我们可以把 $s_2$、$s_1$ 展开为 $s_2 = \alpha s_+ + \beta s_-$、$s_1 = \gamma s_+ + \delta s_-$；一个简单的计算表明
\begin{equation}
\langle L_{2n} \rangle = - \frac{ \frac{\beta \gamma}{\alpha\delta} \lambda^{2n} - 2 + \frac{\alpha \delta }{\beta \gamma} \lambda^{-2n} }{(\lambda - \lambda^{-1})^2}.
\end{equation}
（我们假定 $s_1, s_2$ 处于一般位置。）然后定义
\begin{equation}
e^{ \theta}:= \sqrt{\frac{\beta \gamma}{\alpha\delta} }, \qquad e^{\varphi}:=\lambda,
\end{equation}
便可将 $\sqrt{\langle L_{2n} \rangle}$ 表为
\begin{align}\label{eq:Trig-Ln}
\begin{split}
\sqrt{\langle L_{2n} \rangle } & = i \frac{\sinh(\theta + n \varphi)}{\sinh\varphi} \\
& = i \left( \frac{\sinh\theta}{\sinh\varphi} T_n(\cosh\varphi) + \cosh\theta \, U_{n-1}(\cosh\varphi)\right),
\end{split}
\end{align}
其中我们引入了切比雪夫多项式 $U_n(\cos\theta) := \frac{\sin(n+1)\theta}{\sin\theta}$ 与 $T_n(\cos\theta) := \cos(n\theta)$。特别地，
\begin{equation}\label{eq:Trig-X-Y}
\frac{1}{\sqrt{Y}} = i \frac{\sinh\theta}{\cosh\varphi}, \qquad \sqrt{X} = i \frac{\sinh(\theta+\varphi)}{\sinh\varphi}.
\end{equation}
反解 \eqref{eq:Trig-X-Y} 得到
\begin{align}\label{eq:solve-coshphi}
 2\cosh \varphi &= \frac{1}{\sqrt{XY}} \left( XY + Y + 1 \right), \\
 2\cosh \theta &= \frac{i}{ \sqrt{X}Y}(XY-Y-1),
\end{align}
于是合起来便得到所要的结果
\begin{equation}
\inprod{L_{2n}} = Y^{-1} \left[T_n(\cosh\varphi) + \frac{XY- Y - 1}{2 \sqrt{X Y}} U_{n-1}(\cosh\varphi) \right]^2.
\end{equation}

% --- 本块公式：F.1--F.10（10 条，其中 F.5/F.6、F.9/F.10 为 align/gather 两行）
%     + G.1--G.8（8 条，其中 G.6/G.7 为 align 两行）；合计 18 条，与原文一致。图 4 幅。 ---
